\section{Introduction}\label{sec:introduction}

The integer factorization problem asks, on input a binary integer $N\ge2$,
for the nondecreasing list of primes whose product is $N$, including
multiplicities. We consider quantum circuits on qubits with classical
basis inputs. A circuit family is polynomial-time uniform if a
deterministic classical algorithm generates the circuit for inputs of each
bit length in time polynomial in that length.

\begin{theorem}\label{thm:main}
There is a polynomial-time uniform family of polynomial-size quantum circuits that outputs the complete prime factorization of every input integer $N\ge2$ with probability one. The number of gates and qubits is bounded by a fixed polynomial in the bit length of $N$. One fixed finite gate set suffices: NOT, CNOT, Toffoli, Hadamard, $\operatorname{diag}(1,i)$, their inverses, and their singly controlled versions.
\end{theorem}

\subsection{Factoring and exact quantum computation}

Shor's quantum algorithms established polynomial-time factorization and
discrete logarithms with bounded error~\cite{Shor}. Their order-based
factoring reduction builds on Miller's work~\cite{Miller,Shor}: suitable
multiplicative orders produce nontrivial square roots of unity and hence
gcd splits. The splitting argument used here is proved directly in
Section~\ref{sec:process}. Verifying a proposed factorization makes
repetition reliable, but a procedure that retries until success need not
have a polynomial worst-case running time. Approximating an arbitrary-angle
rotation to increasingly high precision likewise does not by itself yield
probability-one correctness over a fixed finite gate set.

Exact transforms and algorithms have been developed in several quantum
models. Mosca and Zalka~\cite{MZ} construct exact Fourier transforms using
computed rotations and obtain exact discrete-logarithm algorithms for
known group orders. Their discussion also reports Mosca's exact factoring
algorithm in a generalized model permitting gate parameters obtained
during the computation~\cite[Section~5.2]{MZ}; see also
\cite[Section~1]{Imran}. Imran's exact order-finding algorithm takes a
supplied multiple of the unknown order~\cite[Section~3]{Imran}. A suitable
multiple of polynomial bit length is additional information, which the
factoring problem does not provide initially. These are material model
and input distinctions. Nishimura and Ozawa~\cite{NO} study exact
computation by finitely generated uniform circuit families; their results
also explain why approximate universality cannot replace an explicit gate
analysis in an exact algorithm. Theorem~\ref{thm:main} uses the single
finite gate set stated there and receives only the integer to be factored.

Recent quantum factoring algorithms also address resource costs. Regev's
multidimensional construction~\cite{Regev} initially used a
number-theoretic hypothesis; Pilatte~\cite{Pilatte} proves unconditional
correctness for modified algorithms. Their success guarantees retain a
nonzero failure probability. The present construction addresses certainty
and worst-case polynomial bounds in a specified exact gate model. Its
large polynomial parameters are not intended as practical resource
improvements. It makes no claim of a classical deterministic or randomized
polynomial-time factoring algorithm.

The final step belongs to the amplitude-amplification method originating
in quantum search~\cite{Grover,BHMT}. Exact probability engineering already
appears in Brassard and H\o yer's exact algorithm for Simon's
problem~\cite{BH}, and Mosca and Zalka explicitly use success probability
$1/4$ followed by one ordinary amplification step~\cite[Section~2.1]{MZ}.
General exact amplification may employ phases depending on the success
probability~\cite{Hoyer,BHMT}. Here the preceding construction makes that
probability exactly $1/4$ using dyadic arithmetic and the fixed gates, so
the established pair of sign reflections suffices.

\subsection{Proof strategy and the ancestry of its ingredients}

The proof has two main ingredients beyond the order-based reduction to
factoring. First, it constructs an order trial whose mass on the true
order is a prescribed dyadic rational, meaning a rational with a
power-of-two denominator. For a unit $a$ modulo $m$, its order is the least
positive integer $r$ with $a^r\equiv1\pmod m$. Our trial returns no candidate
or a positive multiple of $r$, and its probability of returning $r$ is an
explicit arithmetic function of $r$. A phase register realizes a
piecewise-linear interpolant of the complex exponential. Its use of
controlled modular shifts is related to the phase-kickback construction
of Cleve, Ekert, Macchiavello, and Mosca~\cite[Section~7]{CEMM}, but the
register here is a four-phase step state rather than a Fourier eigenstate.
Testing its return produces the stated interpolated matrix with its
actual, unnormalized probability.

Lev's four-phase approximation to the order-sampling
transform~\cite[Sections~3--4]{Lev} provides a related approach to avoiding
fine phase rotations. The interpolation and exact
corrections below supply a different selected-branch transform and a
specific success law. The residue-class probabilities admit dyadic
majorants whose sums can be evaluated by a bounded number of polynomial
pieces and exact power sums. Mosca and Zalka use power sums to approximate
a transform's success probability to arbitrary precision, then use a
rotation with a computed angle to obtain an exact
transform~\cite[Sections~4.1--4.2]{MZ}. Here the finite polynomial pieces give
exact dyadic sums. Two small correction branches then make
the required true-order mass exact; neither the true order nor its
factorization is needed to run the trial.

Second, a factoring run also generates the prime factorizations of $p-1$
for its prime factors $p>2$, recursively. This auxiliary occurrence tree
has polynomial size. Its recursive shape is related to Pratt's succinct
primality certificates~\cite{Pratt}, although its role here is different:
prime labels are checked by the deterministic primality test proved in
Appendix~\ref{app:primality}, and the auxiliary factorizations determine
orders and transition probabilities. The appendix adapts the
polynomial-congruence method of Agrawal, Kayal, and Saxena~\cite{AKS}, with
its constants and auxiliary-prime choice established directly.

Once the run has completed and its data have been verified, those data
determine the exact probabilities needed to test earlier outputs. Ordinary
transitions are mixed with uniform guesses of their outputs, and filters
are computed from the completed data. Each transition has the same exact
retained mass. The crucial argument identifies these deferred tests with
a hypothetical sequence of tests using the true data: passing that
sequence itself forces the run to generate the data. Both descriptions
mark the same recorded runs. Thus the analysis introduces neither advice
nor conditioning on eventual completion.

This separates an arithmetic task---making a useful event have a
computable exact mass---from a verification task---recovering the
information needed to evaluate that mass on completed runs. Quantum
rejection sampling~\cite{ORR} is a related use of selective retention and
amplification, with different state-conversion inputs and parameterized
rotations. Our one-transition completion argument is stated independently
in Lemma~\ref{lem:completion}; efficient application also requires
verification of the auxiliary data and the all-good generation property.
A second dyadic correction makes the full preparation's success
probability exactly $1/4$.

Section~\ref{sec:model} specifies exact sampling and coherent classical
control. Section~\ref{sec:order} constructs the prescribed-probability
order trial. Section~\ref{sec:process} builds the factoring process and its
auxiliary data. Section~\ref{sec:exactness} equalizes the transition masses,
identifies the deferred and analytical events, and proves
Theorem~\ref{thm:main} by exact amplification.
